Let \(k\) be a field.
Let \(\operatorname{Ch}(k)\) be the category of unbounded chain complexes of \(k\)-vector spaces.
The derived category \(\mathcal D(k)\) is the localization of \(\operatorname{Ch}(k)\) at the class \(\mathcal W\) of quasi-isomorphisms:
\[
\mathcal D(k) = \operatorname{Ch}(k)[\mathcal W^{-1}].
\]
Here a quasi-isomorphism is a chain map inducing an isomorphism on homology.

Let
\[
H \colon \mathcal D(k) \to \mathcal D(k)
\]
be the endofunctor sending a complex \(V\) to its homology \(H(V)\), regarded as a chain complex with zero differential.

Since \(k\) is a field, every short exact sequence of \(k\)-vector spaces splits.
It follows that every chain complex \(V\) is isomorphic in \(\mathcal D(k)\) to the complex \(H(V)\) with zero differential.
Indeed, after choosing splittings, one can decompose \(V\) as the direct sum of \(H(V)\) and a contractible complex.
The splitting is not canonical at the level of chain complexes.
However, the induced isomorphism in \(\mathcal D(k)\) is canonical.

More precisely, for any two complexes \(V\) and \(W\), homology induces a canonical bijection
\[
\operatorname{Hom}_{\mathcal D(k)}(V,W) \cong \prod_{n \in \mathbb Z} \operatorname{Hom}_k(H(V)_n,H(W)_n).
\]
Thus a morphism in \(\mathcal D(k)\) is completely determined by its induced map on homology.

For each object \(V\), define
\[
\eta_V \colon V \to H(V)
\]
to be the unique morphism in \(\mathcal D(k)\) inducing the identity map
\[
\operatorname{id}_{H(V)} \colon H(V) \to H(V)
\]
on homology.
Similarly, define
\[
\epsilon_V \colon H(V) \to V
\]
to be the unique morphism in \(\mathcal D(k)\) inducing the identity map
\[
\operatorname{id}_{H(V)} \colon H(V) \to H(V)
\]
on homology.

The composites
\[
H(V) \xrightarrow{\epsilon_V} V \xrightarrow{\eta_V} H(V)
\]
and
\[
V \xrightarrow{\eta_V} H(V) \xrightarrow{\epsilon_V} V
\]
induce the identity maps on homology.
Therefore both composites are identity morphisms in \(\mathcal D(k)\).
Hence \(\eta_V\) and \(\epsilon_V\) are inverse isomorphisms.

The morphisms \(\eta_V\) are natural in \(V\).
Indeed, for any morphism \(f \colon V \to W\) in \(\mathcal D(k)\), the two composites
\[
V \xrightarrow{\eta_V} H(V) \xrightarrow{H(f)} H(W)
\]
and
\[
V \xrightarrow{f} W \xrightarrow{\eta_W} H(W)
\]
induce the same map
\[
H(f) \colon H(V) \to H(W)
\]
on homology.
Since morphisms in \(\mathcal D(k)\) are detected by homology, the two composites are equal.
Thus the maps \(\eta_V\) assemble into a natural isomorphism
\[
\eta \colon \operatorname{Id}_{\mathcal D(k)} \Rightarrow H.
\]
Similarly, the maps \(\epsilon_V\) assemble into a natural isomorphism
\[
\epsilon \colon H \Rightarrow \operatorname{Id}_{\mathcal D(k)}.
\]
Consequently,
\[
\operatorname{Id}_{\mathcal D(k)} \cong H
\]
as endofunctors of \(\mathcal D(k)\).

This is a statement in the derived category.
It does not assert the existence of canonical chain maps \(V \to H(V)\) or \(H(V) \to V\) in \(\operatorname{Ch}(k)\).
Such chain-level maps require choices of cycles representing homology classes, or equivalently choices of splittings of cycles and boundaries.